\begin{lemma}
Assume \(\noderef{ParameterHierarchy}(\varepsilon,\varepsilon_1,\varepsilon_2,n_0)\).
For every \(\eta>0\) and \(\delta>0\), there is \(N\ge n_0\) such that, for
all \(n>N\), with
\[
K=\noderef{TwoBiteNaturalIndependenceScale}(1+\varepsilon,n),
\]
we have
\[
1\le K\le n,
\qquad
\log\binom nK\le \left(\frac12+\eta\right)K\log n,
\]
and
\[
2\log(K+1)+\log(\delta^{-1})\le \eta K\log n.
\]
\end{lemma}
\begin{proof}
From \(\noderef{ParameterHierarchyEventualComparisons}\), for every
\(n>n_0\),
\[
(1+\varepsilon)\sqrt{n\log n}\le K
<(1+\varepsilon)\sqrt{n\log n}+1.
\]
The hierarchy also gives \(0<\varepsilon<1\).  Put
\[
S=\sqrt{n\log n},\qquad
C_\delta=2\log 4+\max\{0,\log(\delta^{-1})\}.
\]
Choose \(N\ge n_0\) so that for every \(n>N\) the following finite list of
inequalities holds:
\[
\log n\ge 1,\qquad S\ge1,\qquad 9\log n\le n, \tag{1}
\]
\[
\log n\ge \eta^{-1},\qquad \log n\ge C_\delta,\qquad
3\le \eta\sqrt{n\log n}. \tag{2}
\]
Equivalently, it is enough for \(N\) to force the concrete requirements
\[
n>e,\qquad \frac{n}{\log n}\ge9,\qquad
\log n\ge\eta^{-1},\qquad \log n\ge C_\delta,\qquad
n\log n\ge 9\eta^{-2}
\]
for all \(n>N\).  These requirements imply (1) and (2) term by term.
These are the only large-\(n\) margins used below.  With \(S\ge1\) and
\(\varepsilon<1\), the rounded comparison gives
\[
K\ge (1+\varepsilon)S\ge S\ge1. \tag{3}
\]
It also gives
\[
K<(1+\varepsilon)S+1<2S+1\le3S\le n, \tag{4}
\]
where the last inequality is exactly \(9\log n\le n\) after squaring the
nonnegative quantities \(3S\) and \(n\).  Since \(K\) and \(n\) are natural
numbers, (3) and (4) give \(1\le K\le n\).

We next prove the binomial logarithm bound.  For \(1\le K\le n\),
\[
\binom nK
=\frac{n(n-1)\cdots(n-K+1)}{K!}
\le \frac{n^K}{K!}. \tag{5}
\]
The exponential-series estimate
\[
e^K=\sum_{j=0}^{\infty}\frac{K^j}{j!}
\ge \frac{K^K}{K!}
\]
implies
\[
K!\ge \left(\frac K e\right)^K. \tag{6}
\]
Combining (5) and (6) gives
\[
\binom nK\le \left(\frac{en}{K}\right)^K. \tag{7}
\]
Because \(1\le K\le n\), the quantity \(en/K\) is at least \(e>0\), so applying
\(\log\) to (7) gives
\[
\log\binom nK
\le K\log\frac{en}{K}
=K(1+\log n-\log K). \tag{8}
\]
From \(K\ge S=\sqrt{n\log n}\) and \(\log n>0\),
\[
\log K\ge \log S
=\frac12\log n+\frac12\log\log n. \tag{9}
\]
Substituting (9) into (8) gives
\[
\log\binom nK
\le K\left(1+\frac12\log n-\frac12\log\log n\right). \tag{10}
\]
The first inequality in (2) gives \(1\le\eta\log n\), and
\(\log n\ge1\) gives \(\log\log n\ge0\).  Hence
\[
1+\frac12\log n-\frac12\log\log n
\le \left(\frac12+\eta\right)\log n. \tag{11}
\]
Multiplying (11) by the nonnegative integer \(K\) and using (10) proves
\[
\log\binom nK\le\left(\frac12+\eta\right)K\log n.
\]

It remains to prove the finite cell-count and \(\delta\)-margin.  From (4) and
\(S\ge1\),
\[
K+1<2S+2\le4S. \tag{12}
\]
Therefore
\[
\begin{aligned}
2\log(K+1)+\log(\delta^{-1})
&\le 2\log(4S)+\max\{0,\log(\delta^{-1})\}\\
&=C_\delta+\log n+\log\log n.
\end{aligned} \tag{13}
\]
Since \(\log n\ge1\), the elementary inequality \(\log y\le y\) for
\(y\ge1\), applied to \(y=\log n\), gives \(\log\log n\le\log n\).  Together
with \(\log n\ge C_\delta\), this gives
\[
C_\delta+\log n+\log\log n\le 3\log n. \tag{14}
\]
The final condition in (2) gives
\[
3\log n\le \eta\sqrt{n\log n}\,\log n
=\eta S\log n. \tag{15}
\]
Finally, \(K\ge S\) from (3) and \(\eta>0\), \(\log n\ge0\), imply
\[
\eta S\log n\le \eta K\log n. \tag{16}
\]
Combining (13), (14), (15), and (16) gives
\[
2\log(K+1)+\log(\delta^{-1})\le \eta K\log n.
\]
\end{proof}
